%%
%% Guarda o SUMARIO: a grafia de cada um dos cinco niveis, o alinhamento dos
%% indicativos e dos titulos, e o que acontece com um titulo que ocupa duas
%% linhas.
%%
%% A 2.6 manda que "no sumario, as secoes devem ser grafadas conforme
%% apresentadas no corpo do trabalho", e a 3.1.2.1.6 manda alinhar os titulos
%% "pela margem do titulo do indicativo mais extenso" -- e usar como exemplo o
%% sumario do proprio manual, cujo layout e: todo indicativo na margem
%% esquerda, todo titulo numa unica coluna.
%%
%% Cada nivel aparece duas vezes: uma com titulo curto, para medir a coluna, e
%% uma com titulo longo, para ver onde a segunda linha comeca.
%%
%% NAO e gerado pelo coppe.dtx e NAO e distribuido: a suite prova que a
%% classe funciona, nao faz parte dela. Rode por tools\prova.ps1, ou por
%% tests\run-tests.ps1 para so esta pasta.
%%

\documentclass[dsc]{coppe}

\usepackage[brazilian]{babel}

\title{Título de Teste do Sumário}
\foreigntitle{Table of Contents Test}
\author{Nome}{de Teste}

\advisor[UFRJ]{Prof.}{Primeiro}{Orientador}{D.Sc.}
\examiner[UFRJ]{Prof.}{Primeiro Examinador}{D.Sc.}
\examiner{Prof.}{Segundo Examinador}{Ph.D.}

\department{PESC}
\date{09}{2026}
\keyword{Sumário}

\begin{document}
  \maketitle
  \frontmatter
  \begin{abstract}
    Resumo de teste.
  \end{abstract}
  \begin{foreignabstract}
    Test abstract.
  \end{foreignabstract}
  \tableofcontents
  \mainmatter

  \chapter{Curto}
    Texto.
  \section{Curto}
    Texto.
  \subsection{Curto}
    Texto.
  \subsubsection{Curto}
    Texto.
  \paragraph{Curto}
    Texto.

  \chapter{Um título de seção primária deliberadamente longo, escrito para
    ocupar mais de uma linha no sumário e mostrar onde a segunda linha começa}
    Texto.
  \section{Um título de seção secundária deliberadamente longo, escrito para
    ocupar mais de uma linha no sumário e mostrar onde a segunda linha começa}
    Texto.
  \subsection{Um título de seção terciária deliberadamente longo, escrito para
    ocupar mais de uma linha no sumário e mostrar onde a segunda linha começa}
    Texto.
  \subsubsection{Um título de seção quaternária deliberadamente longo, escrito
    para ocupar mais de uma linha no sumário e mostrar onde a segunda começa}
    Texto.
  \paragraph{Um título de seção quinária deliberadamente longo, escrito para
    ocupar mais de uma linha no sumário e mostrar onde a segunda linha começa}
    Texto.
% Indicativos de dois algarismos em todos os niveis. A caixa que o
% \numberline reserva para o indicativo tem largura fixa por nivel; um
% indicativo mais largo que ela nao empurra o titulo, ele entra por cima.
  \setcounter{chapter}{9}
  \chapter{Depois de nove}
    Texto.
  \setcounter{section}{9}
  \section{Depois de nove}
    Texto.
  \setcounter{subsection}{9}
  \subsection{Depois de nove}
    Texto.
  \setcounter{subsubsection}{9}
  \subsubsection{Depois de nove}
    Texto.
  \setcounter{paragraph}{9}
  \paragraph{Depois de nove}
    Texto.
\end{document}
%%
%%
%% End of file `tests/test_sumario.tex'.
